\documentclass[11pt]{article}

\usepackage[margin=1in]{geometry}
\usepackage{amsmath, amssymb}
\usepackage{graphicx}
\usepackage{booktabs}
\usepackage{listings}
\usepackage{xcolor}
\usepackage{float}

\lstset{
    basicstyle=\ttfamily\small,
    breaklines=true,
    frame=single,
    language=Python
}

\title{IE4497 Homework 2}
\author{Wang Zhoufu}
\date{Mar 29 2026}

\begin{document}
\sloppy

\maketitle

\section{Question 1}

Consider a step function $f(x) = 1$ if $x>0$, and $f(x) = 0$ otherwise. Is this a good choice for hidden neuron activation functions when you need to train the network using backprop? Justify your answer.

\textbf{Answer:} 

The step function is not a suitable activation function for training neural networks using backpropagation.

\begin{itemize}
    \item First, the step function is non-differentiable at $x=0$ and has zero derivative everywhere else. Since backpropagation relies on gradients to update weights, this causes a problem: the gradient is either undefined or zero, so no meaningful weight updates can occur. This effectively prevents learning.

    \item Second, because the gradient is zero almost everywhere, the network suffers from the \textbf{vanishing gradient problem} in an extreme form, where learning completely stalls after initialization.

    \item Third, the step function produces binary outputs, which removes sensitivity to small changes in input. This makes optimization difficult, as the network cannot adjust outputs smoothly in response to parameter updates.
\end{itemize}

Therefore, smooth and differentiable activation functions such as $\tanh$, sigmoid, or ReLU are preferred for training neural networks using backpropagation.

\section{Function Analysis}


The target function is
\begin{equation}
g(p) = 1 + \sin\left(\frac{3\pi}{2}p\right).
\end{equation} 

So \begin{equation}
T = \frac{2\pi}{\omega}.
\end{equation}

Therefore, the network is trained and evaluated on the interval
\begin{equation}
p \in \left[0,\frac{4}{3}\right].
\end{equation}

This is a smooth periodic function, which makes it suitable for approximation by a small feedforward neural network.

\section{Data Preprocessing}

\subsection{Dataset Construction}

The dataset is generated directly from the analytical function using uniformly spaced points over one full period:
\begin{equation}
p_{\text{raw}} = \text{linspace}(0,T,\texttt{num\_points}).
\end{equation}
For each sample point $p$, the target value is computed as
\begin{equation}
y = g(p) = 1 + \sin\left(\frac{3\pi}{2}p\right).
\end{equation}

\subsection{Input Normalization}

The raw input $p \in [0,T]$ is normalized to $[-1,1]$ before being passed into the neural network:
\begin{equation}
p_{\text{norm}} = 2\frac{p-0}{T-0} - 1.
\end{equation}

This normalization improves numerical stability during training and helps optimization converge more smoothly.

\subsection{Train--Validation Split}

The code uses a deterministic split:
\begin{equation}
\texttt{split\_idx} = \lfloor \texttt{train\_ratio} \cdot n \rfloor,
\end{equation}
where the first portion of the ordered samples is used for training and the remaining portion is used for validation.

Thus, the split is not random. This detail is important because it means the validation set may occupy a different region of the input interval from the training set, especially in the overfitting experiment.

\section{Model Architecture}

The Python code implements a single-hidden-layer 1--N--1 neural network. The architecture is:

\begin{itemize}
    \item 1 input neuron,
    \item $N$ hidden neurons,
    \item 1 output neuron.
\end{itemize}

The implemented model is
\begin{lstlisting}
self.net = nn.Sequential(
    nn.Linear(1, hidden_dim),
    act_layer,
    nn.Linear(hidden_dim, 1),
)
\end{lstlisting}

\section{Training Methodology}

\subsection{Loss Function}

The network is trained using mean squared error (MSE):
\begin{equation}
\mathcal{L}_{\text{MSE}} = \frac{1}{n}\sum_{i=1}^{n}\left(\hat{y}_i - y_i\right)^2.
\end{equation}

\subsection{Optimizer}

The optimizer is Adam with learning rate
\begin{equation}
\eta = 10^{-3}.
\end{equation}

\subsection{Training Procedure}

For each epoch:
\begin{enumerate}
    \item the model performs a forward pass on the training set;
    \item the training loss is computed;
    \item gradients are backpropagated;
    \item parameters are updated by Adam;
    \item validation loss is computed in evaluation mode.
\end{enumerate}

The code stores both training loss and validation loss for later visualization.

\section{Hyperparameter Search Strategy}

The assignment asks for an explanation of how suitable hyperparameters are found. In this implementation, the study focuses on a small set of interpretable design choices:

\begin{itemize}
    \item hidden width $N$;
    \item number of training epochs;
    \item learning rate;
    \item amount of training data available;
    \item activation function.
\end{itemize}

The final code supports three activation functions:
\begin{itemize}
    \item \texttt{Tanh},
    \item \texttt{ReLU},
    \item \texttt{Sigmoid}.
\end{itemize}

However, the chosen experiments all use \texttt{Tanh}, since it is smooth and naturally suited to approximating a sinusoidal target.

The main hyperparameter logic used in the experiments is:
\begin{itemize}
    \item \textbf{Underfitting:} use lower model capacity and fewer effective training updates;
    \item \textbf{Good fit:} use moderate-to-large capacity and enough epochs;
    \item \textbf{Overfitting:} use very high capacity and a small training portion.
\end{itemize}

\section{Experimental Configurations}

The code runs three experiments through the \texttt{ExperimentConfig} dataclass.

\subsection{Underfitting Configuration}

\begin{table}[H]
\centering
\begin{tabular}{ll}
\toprule
Parameter & Value \\
\midrule
Hidden dimension $N$ & 16 \\
Epochs & 500 \\
Learning rate & $10^{-3}$ \\
Number of points & 512 \\
Training ratio & 0.8 \\
Activation & Tanh \\
\bottomrule
\end{tabular}
\caption{Underfitting experiment configuration.}
\end{table}

\subsection{Good Fit Configuration}

\begin{table}[H]
\centering
\begin{tabular}{ll}
\toprule
Parameter & Value \\
\midrule
Hidden dimension $N$ & 64 \\
Epochs & 5000 \\
Learning rate & $10^{-3}$ \\
Number of points & 512 \\
Training ratio & 0.8 \\
Activation & Tanh \\
\bottomrule
\end{tabular}
\caption{Good fit experiment configuration.}
\end{table}

\subsection{Overfitting Configuration}

\begin{table}[H]
\centering
\begin{tabular}{ll}
\toprule
Parameter & Value \\
\midrule
Hidden dimension $N$ & 128 \\
Epochs & 5000 \\
Learning rate & $10^{-3}$ \\
Number of points & 512 \\
Training ratio & 0.2 \\
Activation & Tanh \\
\bottomrule
\end{tabular}
\caption{Overfitting experiment configuration.}
\end{table}

A key point is that the overfitting case in the code does \textbf{not} reduce the total number of generated points. Instead, it keeps 512 total points but trains on only 20\% of them. Because the split is sequential rather than random, the training and validation inputs come from different subintervals of the domain. Therefore, the overfitting behavior should be interpreted carefully.

\section{Evaluation Metrics}

The code computes three regression metrics:

\begin{align}
\text{MSE} &= \frac{1}{n}\sum_{i=1}^{n}(\hat{y}_i-y_i)^2, \\
\text{RMSE} &= \sqrt{\text{MSE}}, \\
\text{MAE} &= \frac{1}{n}\sum_{i=1}^{n}|\hat{y}_i-y_i|.
\end{align}

These metrics are evaluated separately on the training set and the validation set.

\section{Results and Interpretation}

\subsection{Underfitting}

\begin{figure}[H]
\centering
\includegraphics[width=0.72\textwidth]{plots/underfitting/fit.png}
\caption{Underfitting: network output versus true function.}
\end{figure}

\begin{figure}[H]
\centering
\includegraphics[width=0.72\textwidth]{plots/underfitting/loss.png}
\caption{Underfitting: training and validation loss curves.}
\end{figure}

In the underfitting case, the network has limited representational power relative to the task and training budget. As a result:
\begin{itemize}
    \item the predicted curve does not fully capture the sinusoidal shape;
    \item both training and validation losses remain relatively high;
    \item the model shows high bias.
\end{itemize}

\subsection{Good Fit}

\begin{figure}[H]
\centering
\includegraphics[width=0.72\textwidth]{plots/good_fit/fit.png}
\caption{Good fit: network output versus true function.}
\end{figure}

\begin{figure}[H]
\centering
\includegraphics[width=0.72\textwidth]{plots/good_fit/loss.png}
\caption{Good fit: training and validation loss curves.}
\end{figure}

In the good fit case, the network approximates the target function closely across the full domain. The loss curves show stable convergence, and the prediction follows the true sinusoid with small error. This indicates that the chosen hidden width and training duration are sufficient for this approximation problem.

\subsection{Overfitting}

\begin{figure}[H]
\centering
\includegraphics[width=0.72\textwidth]{plots/overfitting/fit.png}
\caption{Overfitting experiment: network output versus true function.}
\end{figure}

\begin{figure}[H]
\centering
\includegraphics[width=0.72\textwidth]{plots/overfitting/loss.png}
\caption{Overfitting experiment: training and validation loss curves.}
\end{figure}

The code labels the third experiment as ``Overfitting.'' The model has very high capacity and is trained for many epochs using only 20\% of the ordered dataset for training. This setup can produce very low training loss while causing much poorer validation performance.

However, the interpretation requires care:
\begin{itemize}
    \item the training set and validation set are not randomly sampled;
    \item the training data comes from the earlier part of the interval and the validation data from the later part;
    \item therefore, the gap between training and validation performance reflects not only overfitting, but also distribution shift caused by the split.
\end{itemize}

So, in this code, the third case demonstrates a strong \emph{train--validation mismatch}, which is consistent with overfitting, but it is not a perfectly pure overfitting experiment in the statistical sense.

\appendix
\section{Python Source Code}

\begin{lstlisting}[language=Python]
import os
from typing import Literal
from dataclasses import dataclass

import matplotlib.pyplot as plt
import torch
import torch.nn as nn


torch.manual_seed(42)


@dataclass
class ExperimentConfig:
    name: str
    hidden_dim: int
    epochs: int
    lr: float
    num_points: int
    training_ratio: float = 0.8
    activation: Literal["tanh", "relu", "sigmoid"] = "tanh"


class SurrogateNet(nn.Module):
    """
    1-N-1 neural network:
    1 input neuron, N hidden neurons, 1 output neuron.
    """

    def __init__(self, hidden_dim: int, activation: str = "tanh") -> None:
        super().__init__()
        if activation == "tanh":
            act_layer = nn.Tanh()
        elif activation == "relu":
            act_layer = nn.ReLU()
        elif activation == "sigmoid":
            act_layer = nn.Sigmoid()
        else:
            raise ValueError(f"Unsupported activation: {activation}")
        self.net = nn.Sequential(
            nn.Linear(1, hidden_dim),
            act_layer,
            nn.Linear(hidden_dim, 1),
        )

    def forward(self, x: torch.Tensor) -> torch.Tensor:
        return self.net(x)


def target_function(p: torch.Tensor) -> torch.Tensor:
    """
    g(p) = 1 + sin((3*pi/2) * p)
    Period T = 4/3
    """
    return 1 + torch.sin((3 * torch.pi / 2) * p)


def normalize_to_minus1_1(x: torch.Tensor, x_min: float, x_max: float) -> torch.Tensor:
    """
    Normalize x from [x_min, x_max] to [-1, 1].
    """
    return 2.0 * (x - x_min) / (x_max - x_min) - 1.0


def compute_metrics(y_true: torch.Tensor, y_pred: torch.Tensor) -> dict:
    mse = torch.mean((y_pred - y_true) ** 2).item()
    rmse = mse**0.5
    mae = torch.mean(torch.abs(y_pred - y_true)).item()
    return {
        "mse": mse,
        "rmse": rmse,
        "mae": mae,
    }


def make_dataset(
    num_points: int,
    T: float,
) -> tuple[torch.Tensor, torch.Tensor, torch.Tensor]:
    """
    Returns:
        p_raw: raw input in [0, T]
        p_norm: normalized input in [-1, 1]
        y: target values
    """
    p_raw = torch.linspace(0.0, T, num_points).unsqueeze(1)
    y = target_function(p_raw)
    p_norm = normalize_to_minus1_1(p_raw, 0.0, T)
    return p_raw, p_norm, y


def split_dataset(
    x: torch.Tensor,
    y: torch.Tensor,
    train_ratio: float = 0.8,
) -> tuple[torch.Tensor, torch.Tensor, torch.Tensor, torch.Tensor]:
    split_idx = int(train_ratio * len(x))
    x_train, x_val = x[:split_idx], x[split_idx:]
    y_train, y_val = y[:split_idx], y[split_idx:]
    return x_train, y_train, x_val, y_val


def train_model(
    model: nn.Module,
    x_train: torch.Tensor,
    y_train: torch.Tensor,
    x_val: torch.Tensor,
    y_val: torch.Tensor,
    epochs: int,
    lr: float,
) -> tuple[list[float], list[float]]:
    optimizer = torch.optim.Adam(model.parameters(), lr=lr)
    loss_fn = nn.MSELoss()

    train_losses: list[float] = []
    val_losses: list[float] = []

    for epoch in range(epochs):
        model.train()
        optimizer.zero_grad()

        y_pred_train = model(x_train)
        train_loss = loss_fn(y_pred_train, y_train)
        train_loss.backward()
        optimizer.step()

        model.eval()
        with torch.no_grad():
            y_pred_val = model(x_val)
            val_loss = loss_fn(y_pred_val, y_val)

        train_losses.append(train_loss.item())
        val_losses.append(val_loss.item())

        if epoch % 200 == 0 or epoch == epochs - 1:
            print(
                f"Epoch {epoch:4d} | "
                f"Train Loss: {train_loss.item():.6f} | "
                f"Val Loss: {val_loss.item():.6f}"
            )

    return train_losses, val_losses


def plot_fit(
    save_dir: str,
    title: str,
    p_data: torch.Tensor,
    y_data: torch.Tensor,
    p_plot: torch.Tensor,
    y_true_plot: torch.Tensor,
    y_pred_plot: torch.Tensor,
) -> None:
    plt.figure(figsize=(8, 5))
    plt.scatter(
        p_data.numpy(),
        y_data.numpy(),
        s=30,
        alpha=0.7,
        label="Training/validation data",
    )
    plt.plot(p_plot.numpy(), y_true_plot.numpy(), label="True function")
    plt.plot(p_plot.numpy(), y_pred_plot.numpy(), label="Network output")
    plt.xlabel("p")
    plt.ylabel("g(p)")
    plt.title(title)
    plt.legend()
    plt.tight_layout()
    plt.savefig(os.path.join(save_dir, "fit.png"), dpi=200)
    plt.close()


def plot_losses(
    save_dir: str,
    title: str,
    train_losses: list[float],
    val_losses: list[float],
) -> None:
    plt.figure(figsize=(8, 5))
    plt.plot(train_losses, label="Train Loss")
    plt.plot(val_losses, label="Validation Loss")
    plt.xlabel("Epoch")
    plt.ylabel("MSE Loss")
    plt.title(title)
    plt.legend()
    plt.tight_layout()
    plt.savefig(os.path.join(save_dir, "loss.png"), dpi=200)
    plt.close()


def run_experiment(config: ExperimentConfig, T: float = 4.0 / 3.0) -> dict:
    print(f"\n===== {config.name} =====")

    save_dir = os.path.join("plots", config.name.lower().replace(" ", "_"))
    os.makedirs(save_dir, exist_ok=True)

    p_raw, p_norm, y = make_dataset(
        num_points=config.num_points,
        T=T,
    )

    x_train, y_train, x_val, y_val = split_dataset(p_norm, y, config.training_ratio)

    model = SurrogateNet(hidden_dim=config.hidden_dim, activation=config.activation)

    train_losses, val_losses = train_model(
        model=model,
        x_train=x_train,
        y_train=y_train,
        x_val=x_val,
        y_val=y_val,
        epochs=config.epochs,
        lr=config.lr,
    )

    model.eval()
    with torch.no_grad():
        y_train_pred = model(x_train)
        y_val_pred = model(x_val)

    train_metrics = compute_metrics(y_train, y_train_pred)
    val_metrics = compute_metrics(y_val, y_val_pred)

    print("Train Metrics:")
    print(
        f"  MSE={train_metrics['mse']:.6f}, "
        f"RMSE={train_metrics['rmse']:.6f}, "
        f"MAE={train_metrics['mae']:.6f}"
    )
    print("Validation Metrics:")
    print(
        f"  MSE={val_metrics['mse']:.6f}, "
        f"RMSE={val_metrics['rmse']:.6f}, "
        f"MAE={val_metrics['mae']:.6f}"
    )

    p_plot = torch.linspace(0.0, T, 400).unsqueeze(1)
    p_plot_norm = normalize_to_minus1_1(p_plot, 0.0, T)
    y_true_plot = target_function(p_plot)

    with torch.no_grad():
        y_pred_plot = model(p_plot_norm)

    plot_fit(
        save_dir=save_dir,
        title=f"{config.name}: Network Output vs True Function",
        p_data=p_raw,
        y_data=y,
        p_plot=p_plot,
        y_true_plot=y_true_plot,
        y_pred_plot=y_pred_plot,
    )

    plot_losses(
        save_dir=save_dir,
        title=f"{config.name}: Training and Validation Loss",
        train_losses=train_losses,
        val_losses=val_losses,
    )

    return {
        "config": config,
        "train_metrics": train_metrics,
        "val_metrics": val_metrics,
    }


def main() -> None:
    """
    Recommended setup for the assignment:

    1. Underfitting:
       - very small hidden layer
       - limited epochs
       - full dense dataset

    2. Good Fit:
       - moderate hidden layer
       - enough epochs
       - full dense dataset

    3. Overfitting:
       - large hidden layer
       - many epochs
       - very small dataset
       This makes overfitting actually visible.
    """
    experiments = [
        ExperimentConfig(
            name="Underfitting",
            hidden_dim=16,
            epochs=500,
            lr=1e-3,
            num_points=512,
            activation="tanh",
            training_ratio=0.8,
        ),
        ExperimentConfig(
            name="Good Fit",
            hidden_dim=64,
            epochs=5000,
            lr=1e-3,
            num_points=512,
            activation="tanh",
            training_ratio=0.8,
        ),
        ExperimentConfig(
            name="Overfitting",
            hidden_dim=128,
            epochs=5000,
            lr=1e-3,
            num_points=512,
            activation="tanh",
            training_ratio=0.2,
        ),
    ]

    results = []
    for config in experiments:
        result = run_experiment(config)
        results.append(result)

    print("\n===== Summary =====")
    for result in results:
        cfg = result["config"]
        val = result["val_metrics"]
        print(
            f"{cfg.name:12s} | "
            f"N={cfg.hidden_dim:<3d} | "
            f"points={cfg.num_points:<3d} | "
            f"Val MSE={val['mse']:.6f} | "
            f"Val RMSE={val['rmse']:.6f} | "
            f"Val MAE={val['mae']:.6f}"
        )


if __name__ == "__main__":
    main()

\end{lstlisting}
\end{document}